% Fichier LaTeX enrichi avec images UPDF
% Exercice 2 - Amérique du Nord ^
% 23 points

\documentclass[12pt]{article}
\usepackage[utf8]{inputenc}
\usepackage[T1]{fontenc}
\usepackage[french]{babel}
\usepackage{graphicx}
\usepackage{amsmath}
\usepackage{amssymb}
\usepackage{geometry}
\geometry{a4paper, margin=2cm}

\begin{document}

\section*{Exercice 2 (23 points)}

% Images de l'exercice
\begin{center}
\includegraphics[width=0.8\textwidth]{images/dnb_2025_06_ameriquenord_2_Image_005.gif}
\includegraphics[width=0.8\textwidth]{images/dnb_2025_06_ameriquenord_2_Image_006.gif}
\includegraphics[width=0.8\textwidth]{images/dnb_2025_06_ameriquenord_2_Image_033.gif}
\end{center}

\begin{enumerate}
\item Calculer la longueur AB.

\item Montrer que les droites (DE) et (BC) sont parallèles.

\item Montrer que la longueur DE est égale à $42$~m.

\item Montrer que la longueur EM est environ égale à $24,2$~m.

\item En déduire l'aire du triangle AMD.

\end{enumerate}

\end{document}
